%% Supplement tables relocated from the body to mirror the canonical PAPER_SUBMISSION.md
%% "## Supplementary Material" structure (S.4 = screening log Table S1; S.5 = titration
%% ladder Table S2). The tables themselves are non-float, \textbf-captioned blocks emitted
%% by _bestrec_run/emit_latex_tables.py from the md; this file only supplies the S.4/S.5
%% section frames and \input positions. See BUILD_NOTES.md.
\section*{S.4 --- Screening-log detail (full probe table)}

\begin{quote}
\emph{Supplement to \S\ref{sec:5.5} --- the full search record. Each row's \ensuremath{\Delta} is versus its stated base; most rows are single-seed exploratory probes (power note in \S\ref{sec:5.5}).}
\end{quote}

%% GENERATED by _bestrec_run/emit_latex_tables.py -- DO NOT EDIT BY HAND
% table key: table2
% provenance: md lines 440-461; family=hstu_tables.json:table2; check={"checked": 33, "exact": 17, "within_1ulp": 13, "allowlisted": 3, "fail": []}
\par\medskip\begingroup\scriptsize\setlength{\tabcolsep}{4pt}\renewcommand{\arraystretch}{1.12}
\noindent \textbf{Table 2: Systematic negative-result map --- capacity-adding / alternative-mechanism probes, all neutral or harmful.} (Mechanism class: where the lever acts. ``Learned scalar \ensuremath{\rightarrow} 0?'' = did a free learnable gate/temperature/intensity collapse the mechanism. n = seed count. \textbf{Note: probes are not equally powered} --- most rows are single-seed exploratory probes; only the multi-seed pre-declared negatives carry confirmatory weight, per the paragraph above.)\par\smallskip
\noindent
\begin{tabularx}{\linewidth}{p{0.155\linewidth}p{0.13\linewidth}p{0.11\linewidth}p{0.16\linewidth}p{0.10\linewidth}Y}
\toprule
Lever & Mechanism class & Base (NDCG@10) & \ensuremath{\Delta} NDCG@10 (n) & Learned scalar \ensuremath{\rightarrow} 0? & Verdict \\
\midrule
c3 continuous time-decay attention kernel & attention re-weighting (capacity-add) & SBERT stack \ensuremath{\approx}0.0639 & \textbf{\ensuremath{-}0.0150} (1; test 0.0489 at record val 0.0725) & n/a & Rejected --- actively harmful; the largest validation--test divergence in the corpus \\
Sampled softmax (K=512 negatives) & loss approximation & full-softmax stack & \ensuremath{-}0.0026 (1) & n/a & Rejected --- undertrains full-catalog ranking \\
Dual text encoder (SBERT \ensuremath{\oplus} BLaIR) & text encoder & SBERT stack & \ensuremath{-}0.0014 (1) & n/a & Rejected --- the text encoder is not the gap \\
BLaIR text encoder (swap) & text encoder & SBERT stack & \ensuremath{-}0.0013 (1) & n/a & Rejected --- encoder axis fully swept \\
GD1 spectral-shrink prior & representation regularizer & V2 0.0673 & \ensuremath{-}0.003 (1) & shrink \ensuremath{\rightarrow} 0 & Rejected --- BBP irreducibility figure (\cref{fig:bbp_irreducibility}) retained \\
n\_heads = 4 & architecture (capacity-add) & HSTU-style stack & \ensuremath{-}0.0005 (1) & n/a & Rejected --- ties 2-head even at the published dv=dqk=16 spec \\
CL4SRec self-supervision & SSL auxiliary loss & SBERT stack & \ensuremath{-}0.0005 (1) & n/a & Rejected --- early-epoch only \\
c1 prototype-routed expert heads & capacity-add & SBERT stack & \ensuremath{-}0.0004 (1) & n/a & Rejected --- redundant with TAPE \\
c2 text-distillation aux loss & auxiliary loss & SBERT stack & \ensuremath{-}0.0004 (1) & n/a & Rejected --- redundant with the additive text path \\
text-init / text warm-start & embedding initialization & SBERT stack & \ensuremath{-}0.0003 (1) & n/a & Rejected --- redundant with the additive text path \\
EMA / SWA weight averaging & weight-space regularizer & SBERT stack & +0.0001 (1) & n/a & Rejected --- val-high/test-flat; does not repair the structural gap \\
text-sim bias & attention bias & HSTU-style stack & \ensuremath{\pm}0.0001 (4) & n/a & NEUTRAL --- dead weight, drop candidate \\
W1 niche-share fitness-sharing penalty & inter-item loss penalty & U2/LS0.2 0.0649\ensuremath{\pm}0.0002 & 0.0652 (1), within band & \textbf{\ensuremath{\beta} = \ensuremath{-}7.42} (sign-flipped: boosts, not penalizes, crowded niches) & Rejected --- refutes the anti-crowding remedy \\
X1 frequency-adaptive James--Stein shrinkage & representation regularizer & V2 0.0674\ensuremath{\pm}0.0003 & 0.0673 (1), within band & \textbf{c \ensuremath{\approx} 0} (\ensuremath{\lambda}\_max \ensuremath{\approx} 0.004) \ensuremath{\rightarrow} off & Rejected --- model voted shrinkage off \\
Y1 heat-kernel/manifold label smoothing & loss-target reshaping & V2 0.0674\ensuremath{\pm}0.0003 & 0.06653 (1), \ensuremath{\le} band & \textbf{T \ensuremath{\rightarrow} 0.00061} (\textless{} init) \ensuremath{\rightarrow} delta target & Rejected --- text-manifold geometry voted off \\
Z1 forced ID\ensuremath{\rightarrow}text routing & representation routing & V2 & tail \ensuremath{-}75\% (globally dead) & forced gate & Rejected --- globally harmful \\
CF1 cue-fusion gate & representation fusion & V2 / MI & flat/dead (both datasets) & gate \ensuremath{\rightarrow} 0 & Rejected --- new-lever search closed \\
conn-gate connectivity-gated cold-start ID\ensuremath{\leftrightarrow}text fusion & representation fusion/routing & V2 text stack, MI & \textbf{tail} \ensuremath{\Delta} +0.0001 \ensuremath{\pm} 0.0003, paired 95\% CI {[}\ensuremath{-}0.00025, +0.00045{]} (5); overall flat & \textbf{\ensuremath{\alpha} \ensuremath{\rightarrow} 0.0011 \textless{} init 0.0025} (voted off) & Rejected --- operationalizes the intervention-supported connectivity tail mechanism, but is not actionable (CI includes 0) \\
max\_seq\_len 200 / seq \textgreater{} 50 & sequence length & SBERT stack & \textasciitilde0 (1) & n/a & Rejected --- sequences short post-5-core \\
cosine scoring & scoring function & winning op-point & \ensuremath{\le} 0 (1) & n/a & Rejected --- dot-product better at the op-point \\
\bottomrule
\end{tabularx}
\endgroup\par\medskip


\clearpage
\section*{S.5 --- Titration ladder (full six rungs)}

\begin{quote}
\emph{Supplement to \S\ref{sec:5.4}; integrity-gated against its prose values by the released titration-table script.}
\end{quote}

%% GENERATED by _bestrec_run/emit_latex_tables.py -- DO NOT EDIT BY HAND
% table key: table1e
% provenance: md lines 365-372; family=hstu_tables.json:table1e; check={"checked": 61, "exact": 61, "within_1ulp": 0, "allowlisted": 0, "fail": []}
\par\medskip\begingroup\scriptsize\setlength{\tabcolsep}{4pt}\renewcommand{\arraystretch}{1.12}
\noindent \textbf{Table 1e: the interaction-thinning density-titration ladder} (AR2023 Video\_Games 5-core LLOO, full-catalog n\_eval = 94,762, tail\_n = 10,900; paired text\ensuremath{-}ID, best-by-val; 5 seeds = 20260608--12 per rung; mean \ensuremath{\pm} sample-std, positive-seed count in parens). Realized interactions/item are the run-log values; \ensuremath{\alpha} = (interactions/item) / d\_eff with \textbf{d\_eff = 23} (the GD/MP-kept effective rank of the 64-d item table, \S\ref{sec:5.4} BBP analysis), so \ensuremath{\alpha}(\ensuremath{\rho}=0.66) = 0.700 reproduces the MI-subcritical anchor. All cells are re-read from the frozen \texttt{results\_TITR*/TAIL\_*} JSONs by \texttt{\_bestrec\_run/make\_table\_5\_4\_titration.py}, which \textbf{integrity-gates} the recomputed NDCG head/tail means against the locked prose values above (aborts on any \textgreater5e-5 drift); the HR columns and per-rung seed bands are the net-new tabulation the prose summarized only as Spearman trends.\par\smallskip
\noindent
\begin{tabularx}{\linewidth}{lccYYYY}
\toprule
\ensuremath{\rho} & kept inter./item & \ensuremath{\alpha}=ipp/d\_eff & head \ensuremath{\Delta}NDCG@10 & head \ensuremath{\Delta}HR@10 & tail \ensuremath{\Delta}NDCG@10 & tail \ensuremath{\Delta}HR@10 \\
\midrule
1.00 & 24.405 & 1.061 & +0.002213 \ensuremath{\pm} 0.000241 (5/5) & +0.004238 \ensuremath{\pm} 0.000765 (5/5) & \ensuremath{-}0.000148 \ensuremath{\pm} 0.000179 (2/5) & +0.000128 \ensuremath{\pm} 0.000582 (3/5) \\
0.94 & 22.947 & 0.998 & +0.002389 \ensuremath{\pm} 0.000267 (5/5) & +0.004083 \ensuremath{\pm} 0.000830 (5/5) & +0.000165 \ensuremath{\pm} 0.000370 (4/5) & +0.000807 \ensuremath{\pm} 0.000712 (5/5) \\
0.91 & 22.210 & 0.966 & +0.002544 \ensuremath{\pm} 0.000491 (5/5) & +0.004940 \ensuremath{\pm} 0.000892 (5/5) & +0.000039 \ensuremath{\pm} 0.000310 (3/5) & +0.000110 \ensuremath{\pm} 0.000766 (3/5) \\
0.88 & 21.484 & 0.934 & +0.002520 \ensuremath{\pm} 0.000421 (5/5) & +0.004869 \ensuremath{\pm} 0.000765 (5/5) & +0.000537 \ensuremath{\pm} 0.000386 (4/5) & +0.000661 \ensuremath{\pm} 0.000995 (4/5) \\
0.78 & 19.030 & 0.827 & +0.002664 \ensuremath{\pm} 0.000398 (5/5) & +0.004467 \ensuremath{\pm} 0.001082 (5/5) & +0.000056 \ensuremath{\pm} 0.000552 (2/5) & +0.000239 \ensuremath{\pm} 0.001087 (2/5) \\
0.66 & 16.109 & 0.700 & +0.003540 \ensuremath{\pm} 0.000416 (5/5) & +0.005720 \ensuremath{\pm} 0.000904 (5/5) & \ensuremath{-}0.000108 \ensuremath{\pm} 0.000503 (1/5) & +0.000073 \ensuremath{\pm} 0.001165 (3/5) \\
\bottomrule
\end{tabularx}
\endgroup\par\medskip


\section*{S.6 --- Software V3 paired-seed detail}

\begingroup
\renewcommand{\thefigure}{S\arabic{figure}}
\setcounter{figure}{0}
\begin{figure}[h]
  \centering
  \includegraphics[width=0.96\linewidth]{fig_software_v3_pairs.pdf}
  \Description{Two panels visualize the eight Software V3 matched seed blocks. The left panel connects each identity NDCG at 10 value to its learned-FIR value. The right panel shows all eight learned-minus-identity differences, the zero line, the predeclared plus 0.000500 reporting threshold, and the paired mean with its ordinary 95 percent paired-t confidence interval.}
  \caption{Software V3 registered paired-seed detail. Every matched seed block has higher sealed TEST NDCG@10 under learned FIR than identity; the second panel shows all paired differences, the frozen +0.000500 threshold, and the ordinary paired-t mean and 95\% CI. This is outcome-known same-team robustness, not independent confirmation. Exact values are in \texttt{figures/fig\_software\_v3\_pairs\_data.csv}.}
  \label{fig:software_v3_pairs}
\end{figure}

\clearpage
\section*{S.7 --- Selected-checkpoint FIR diagnostics}

\begin{figure}[h]
  \centering
  \includegraphics[width=0.96\linewidth]{fig_fir_response.pdf}
  \Description{Two-panel descriptive FIR diagnostic. Panel A plots the
  signed channel-mean residual tap and mean absolute residual tap across
  causal lags 0 through 15, with uncertainty bands over eight seed summaries.
  Panel B plots the magnitude of the effective residual frequency response
  across normalized angular frequency, with an identity-magnitude reference
  at one and an uncertainty band over the same eight seeds.}
  \caption{Learned residual-tap and effective frequency-response summaries from
  the eight released selected checkpoints of the outcome-known active-control
  learned arm. Channels are summarized within each seed; bands are ordinary
  95\% $t$ intervals over the eight seed summaries. Lag 0 is contemporaneous.
  These are descriptive fitted-operator diagnostics, not a mechanism or
  independent-confirmation test; plotted values are released in
  \texttt{figures/fig\_fir\_response\_data.csv}. The effective kernel is the
  residual identity plus the learned $\Delta$ taps.}
  \label{fig:fir-response}
\end{figure}
\endgroup
